 \documentclass[12pt,english]{article}
    \usepackage[T1]{fontenc}
    \usepackage[ansinew]{inputenc}
%    \usepackage{a4}
    \usepackage{babel}
    \usepackage{graphicx}
    \usepackage{fullpage}
\begin{document}

\appendix
\begin{center}
\center{\bf Appendix} \\

\noindent {\bf Extensions to include feedback from abroad}
\end{center}
In order to deal with repercussions (feedback) from abroad we
formulate a simple two-country income model:
\begin{eqnarray}
 Y_i &=& C_i(Y_i-T_i) + I_i + A_i + Ex_i - Im_i, i=1,2. \label{eq:e0}
% Y_2 &=& C_1(Y_2-T_2) + I_2 + A_2 + Ex_2 - Im_2.  \label{eq:e1}
\end{eqnarray}
Its differential form is given by:
\begin{eqnarray}
 dY_i &=& c_i (dY_i-dT_i)  + dA_i + dI_i +  dEx_i - dIm_i, i=1,2.  \label{eq:e4}
% dY_2 &=& c_2 (dY_2-dT_2)  + dA_2 + dEx_2 - dIm_2.  \label{eq:e5}
\end{eqnarray}

Here, the export function of  one country is the import function of
the other country.
\begin{eqnarray}
Ex_j &=&  Im_i = \gamma_i c_1(Y_i-T_i)+ \alpha_i A_i + \iota_i I_i,
i=1,2; j=1,2, i\neq j. \label{eq:e2}
%Ex_1 &=&  Im_2 = \gamma_2 c_2(Y_2-T_2)+ \alpha_2 A_2 + \iota_2 I_2 .
\label{eq:e3}
\end{eqnarray}

The import functions incorporate the changes argued for in the main
text. Shifts in autonomous public expenditures, autonomous
investment imply shifts in the import function of a country. In
order to simplify the exposition we have formally suppressed
(changes) in autonomous consumption. We shall now also exclude
changes in autonomous investment:
 \begin{eqnarray}
 dI_i& = 0, i=1,2
\end{eqnarray}
Changes of imports and exports are then given by:
\begin{eqnarray}
dEx_j &=&  dIm_i = \gamma_i c_i(dY_i-dT_i)+ \alpha_i dA_i, i =1,2,
j=1,2, i\neq j \label{eq:e6}
% dEx_1 &=& dIm_2 = \gamma_2 c_2(dY_2-dT_2)+ \alpha_2 dA_2.
 \label{eq:e7}
\end{eqnarray}
In addition, we define, first, the marginal propensity to save:
\begin{eqnarray}
s_i & = & 1-c_i, i=1,2
\end{eqnarray}
and, second,
\begin{eqnarray}
 S_i  &\stackrel{def}{=}& [ (1-c_i) + c_i
\gamma_i],\hspace{2mm} i=1,2 \\
& =& (s_i + m_i).
\end{eqnarray}
For the changes in income we obtain:
\begin{eqnarray}
 d Y_i & = & \frac{1}{S_i}\{(1-\alpha_i) dA_i -c_i(1-\gamma_i) dT_i \} \\
 & & + \frac{\gamma_j c_j}{S_i}(d Y_j-dT_j)+
 \frac{\alpha_j}{S_i}dA_j, i=1,2, j=1,2, i \neq j.
% d Y_2 & = & \frac{1}{S_2}\{(1-\alpha_2) dA_2 -c_2(1-\gamma_2) dT_2 \} +
% \frac{\gamma_1 c_1}{S_2}d Y_1.
 \end{eqnarray}

  We now consider the change in income that follows from a change in public expenditures in
  the home country only. From
\begin{eqnarray}
d A_2 &=& d T_2 = 0
 \end{eqnarray}
follows:
\begin{eqnarray}
 dY_2 &= & \frac{\gamma_1 c_1}{S_2 }dY_1.
 \end{eqnarray}

We deal with two cases. First, the
  Haavelmo-case where a tax finance constraint is
  included explicitly. Second, the case where the tax financing restriction is excluded from the analysis.


1. Haavelmo-case

With the tax financing constraint
 \begin{eqnarray}
 dT_1 = dA_1
\end{eqnarray}
we obtain
\begin{eqnarray}
d Y_1 & = & (1- \frac{\alpha_1}{S_1}) dA_1 +  \frac{\gamma_2 c_2}{S_1}d Y_2,  \\
dY_1 &=& \frac{1- \frac{\alpha_1}{S_1}}{1
-\frac{\gamma_1c_1}{S_1}\frac{\gamma_2c_2}{S_2}} dA_1.
\end{eqnarray}
\begin{eqnarray}
 m_i & \stackrel{def}{=} & \gamma_i c_i  \hspace{2mm}
 \end{eqnarray}
\begin{eqnarray}
&=& \frac{1- \frac{\alpha_1}{S_1}}{1
-\frac{m_1}{S_1}\frac{m_2}{S_2}} dA_1
\end{eqnarray}
2. Case without tax financing

With
\begin{eqnarray}
 dT_1 =0
 \end{eqnarray}
we obtain
\begin{eqnarray}
d Y_1 & = & \frac{1-\alpha_1}{S_1} dA_1 +  \frac{\gamma_2 c_2}{S_1}d Y_2  \\
dY_1 &
=&\frac{\frac{1-\alpha_1}{S_1}}{1-\frac{\gamma_1c_1}{S_1}\frac{\gamma_2c_2}{S_2}}
dA_1 \\
& = &
\frac{\frac{1-\alpha_1}{S_1}}{1-\frac{m_1}{S_1}\frac{m_2}{S_2}}
\end{eqnarray}

Thus, we find that feedbacks from abroad do modify the results of
the main text for the income change in the home country. Those
results are multiplied by the factor:
\begin{eqnarray}
\frac{1}{1-\frac{\gamma_1c_1}{S_1}\frac{\gamma_2c_2}{S_2}}.
\end{eqnarray}
With the usual assumptions about the parameters included in the
factor, the factor is larger than 1. Thus the feedback magnifies the
absolute value of an income effect but does not change its sign. A
negative value does remain a negative value. A positive value of the
income effect does remain positive. However it may change from a
positive value smaller than 1 to one larger than 1.
\end{document}
